\subsection{Commutative diagrams}

Let, for example, B, C, D, E, F be five sets, and let \(f\) be a map from E to F, \(g\) be a map from B to C, \(h\) be a map from D to E, \(u\) be a map from B to D, and \(v\) be a map from C to E. To summarize a situation of this kind, one often makes use of diagrams; for example, one summarizes the preceding situation by the following diagram (E, II, p.~14):

\[
\begin{array}{c} \mathrm{B} \xrightarrow {g} \mathrm{C} \\ u \Big \downarrow \\ \mathrm{D} \xrightarrow [ h ]{} \mathrm{E} \xrightarrow [ f ]{} \mathrm{F}. \end{array}\tag{1}
\]

In such a diagram, the group of symbols E \(\xrightarrow{f}\) F schematizes the fact that f is a map from E to F. When there can be no ambiguity about f, one suppresses the letter f and simply writes E \(\rightarrow\) F.

When B, C, D, E, F are groups (resp. A-modules) and f, g, h, u, v are homomorphisms of groups (resp. A-modules), one says, for brevity, that diagram (1) is a diagram of groups (resp. of A-modules).

In principle, a diagram is not a mathematical object, but only a figure, intended to facilitate the reading of a proof. In practice, one often uses diagrams as abbreviating symbols, which avoid naming all the sets and all the maps that one wishes to consider; one thus says “consider diagram (1)” instead of saying: “let B, C, D, E, F be five sets\ldots{} and v a map from C to E”; see for example the statement of prop. 1 of no. 2.

Consider, for example, the following diagram:

\[
\begin{array}{c} \mathrm{B} \xrightarrow {f} \mathrm{C} \xrightarrow {g} \mathrm{D} \xrightarrow {h} \mathrm{E} \\ b \Big \downarrow \qquad \qquad c \Big \downarrow \qquad \qquad d \Big \downarrow \qquad e \Big \downarrow \\ \mathrm{B} ^ {\prime} \xrightarrow {f ^ {\prime}} \mathrm{C} ^ {\prime} \xrightarrow {g ^ {\prime}} \mathrm{D} ^ {\prime} \xrightarrow {h ^ {\prime}} \mathrm{E} ^ {\prime}. \end{array}\tag{2}
\]

To every path composed of a certain number of segments of the diagram traversed in the direction indicated by the arrows, one associates a map from the set represented by the origin of the first segment to the set represented by the endpoint of the last segment, namely the composite of the maps represented by the various segments traversed. For every vertex of the diagram, for example C, one agrees that there is a path reduced to C and one associates with it the identity map 1 \(_{C}\) .

In (2), there are, for example, three paths starting from B and ending at D'; the corresponding maps are \(d \circ g \circ f\) , \(g' \circ c \circ f\) and \(g' \circ f' \circ b\) . One says that a diagram is commutative if, for every pair of paths of the diagram having the same origin and the same endpoint, the two corresponding maps are equal; in particular, if a path has its endpoint coinciding with its origin, the corresponding map must be the identity.

For diagram (2) to be commutative, it is necessary and sufficient that the relations

\[
f ^ {\prime} \circ b = c \circ f, \qquad g ^ {\prime} \circ c = d \circ g, \qquad h ^ {\prime} \circ d = e \circ h;\tag{3}
\]

hold; in other words, it is necessary and sufficient that the three square diagrams extracted from (2) be commutative. Indeed, relations (3) imply \(d \circ g \circ f = g' \circ c \circ f\) since \(d \circ g = g' \circ c\) and \(g' \circ c \circ f = g' \circ f' \circ b\) since \(c \circ f = f' \circ b\) ; therefore the three paths starting from B and ending at D' give the same map. One verifies similarly that the four paths starting from B and ending at E' (resp. the three paths starting from C and ending at E') give the same map. Relations (3) mean that the two paths starting from B (resp. C, D) and ending at C' (resp. D', E') give the same map. All the other pairs of vertices of (2) can be joined by at most one path, and diagram (2) is therefore indeed commutative.

In what follows, we shall leave to the reader the task of formulating and verifying analogous results for other types of diagrams.

